\documentclass[a4paper]{scrartcl}

% !TEX root = ./ausarbeitung.tex

\usepackage{fontspec}
\usepackage{polyglossia}
\setdefaultlanguage{english}
\setmainfont{Libertinus Serif}
\setsansfont{Libertinus Sans}[scale=MatchLowercase]
\setmonofont{Libertinus Mono}[scale=MatchLowercase, FakeStretch=0.8]

% Seitenlayout
\usepackage{scrlayer-scrpage}
\ihead{Fortgeschrittenen-Praktikum F75}    % <-------------------------- Header/Name!
\ohead{Frederik Hollis, Justus Zimmermann}
\setcounter{tocdepth}{2}

% Mathe 
\usepackage{amsmath}
\usepackage{mathtools}
\numberwithin{equation}{section}

\usepackage[
    warnings-off={mathtools-colon,mathtools-overbracket},
    math-style=TeX,
    bold-style=TeX,
    partial=upright,
    nabla=upright
]{unicode-math}
\setmathfont{Libertinus Math}

\usepackage{graphicx} % Graphiken

\usepackage{microtype} % Spacing
\usepackage{pdfpages} % PDFs einfügen
\usepackage{booktabs} % Bessere Tabellen
\usepackage{subcaption} % Bessere Captions
\captionsetup{margin=10pt,font=small,labelfont=sc}
\usepackage{diffcoeff} % Ableitungen
\usepackage{listings}

\usepackage[autostyle=true]{csquotes}    % Language aware quotations
\newcommand{\quoteactual}[1]{\enquote{#1}}   % Quoting someone's words
\newcommand{\quotemention}[1]{\enquote{#1}}  % Talk about the word 'word'
\newcommand{\quoteapprox}[1]{\enquote{#1}}   % Saying something not quite right
\newcommand{\quotedef}[1]{\enquote{#1}}    % Defining something

\usepackage{siunitx} % Einheiten
\sisetup{uncertainty-mode=separate}
\sisetup{exponent-product=\ensuremath{\cdot}}
\sisetup{power-half-as-sqrt}
\sisetup{output-decimal-marker={.}}
\sisetup{table-align-uncertainty=true}
% \sisetup{list-final-separator = {,\,}}    % \qtylist{numbers}{unit}, \numlist{numbers}
% \sisetup{list-separator = {,\,}}    % keep spacing the same between numbers
% for listing multiple quantities if you want compacter list of \num{1;2;3} with better spacing/remove "and"

\DeclareSIUnit{\Vss}{V\textsubscript{SS}}
\DeclareSIUnit{\VRMS}{V\textsubscript{RMS}}
\DeclareSIUnit{\decibelvolt}{dBV}
\DeclareSIUnit{\sigma}{σ}

% Bibliography
\usepackage[
    backend=biber,
    style=phys,
    sorting=none
]{biblatex}
\addbibresource{./ref.bib}

\usepackage[pageanchor]{hyperref} % Links
\definecolor{myNavy}{RGB}{0,0,80}
\hypersetup{
    colorlinks=true,
    linkcolor=myNavy,
    urlcolor=blue,
}
\urlstyle{same}


\title{F75: Neuromorphic Computing}
\author{Justus Zimmermann and Frederik Hollis}
\date{\today}
\publishers{Tutor: Pablo Jimenez}

\begin{document}

\maketitle
\tableofcontents

\include{./...}
A
Nehmen Sie die Resonanzkurven fur die drei Frequenzen ν0, ν1 und ν2 auf.

B
Temperaturabhangigkeit der Resonanzfrequenz

I
Sie nehmen eine komplette Resonanzkurve fur jede eingestellte Temperatur.

II
Mit dieser Methode kann man die Temperaturabhangigkeit viel schneller messen und ohne Nacht- ¨
messung auskommen. Man benutzt dabei das A cos ϕ-Signal des Lock-In-Verstarkers. Zun ¨ achst stellt ¨
man die Spannung im Peltier-Element entsprechend der Temper




(
• Was wollen wir im Abschnitt A und B messen?
• Versuchsplanung: Wie wollen Sie die Resonanzfrequenzen im Abschnitt A finden?
• Was sind Vor- und Nachteile der beiden Methoden fur die Messung im Abschnitt B?
)


Auswertung

Bei der Datenauswertung lernen wir Antworten auf zwei fundamentale Fragen:
• Wie werden die Werte der Parameter eines theoretischen Modells ermittelt, mit denen die
Daten optimal beschrieben werden?
• Wie pr ¨uft man quantitativ, ob das Modell die Daten beschreibt?


“durch Fits mit der χ
2
-Methode”

Least squares

A
Die aufgenommenen Resonanzkurven sollen graphisch dargestellt werden. Genauer gesagt sollen
die Kurven fur die Messgr ¨ oßen ¨ X und Y sowie die abgeleitete Phase gegen die Frequenz aufgetragen werden. Weiterhin sind mit Hilfe eines nichtlinearen parametrischen χ
2
-Fits die Parameter
der Funktion aus Gl. (1.2) an die gemessenen Daten der Amplitude anzupassen

p-value
residuals graph

Gute des Resonators ¨ Q mittels Gl. (1.6) 
Die Gute kann man auch aus der Breite der Kurven gem ¨ aß Gl. (1.8) 
Bestimmen Sie die Gute f ¨ ur mindestens ein Re- ¨
sonanzsignal mit einer diesen Methoden und vergleichen Sie das Ergebnis mit dem aus dem Fit
gewonnenen Wert.

See Anleitung for Fehlerquellen

B
I
Wenn Sie die Methode I benutzt haben, ermitteln Sie jetzt fur jede eingestellte Temperatur mit einem ¨
χ
2
-Fit die Resonanzfrequenz mit ihrem Fehler. Die Resonanzfrequenzen sollten, gegen die Temperatur aufgetragen, naherungsweise auf einer Geraden liegen

II
Bei der Methode II haben Sie die Werte bereits bei der Messung gespeichert. Hier gibt es keine
Abschatzung f ¨ ur die Fehler der Frequenzmessungen. Deswegen k ¨ onnen Sie einen ¨ ublichen ¨ χ
2
-Fit
mit den Fehlern der Temperaturmessung durchfuhren.


Aus der Steigung ergibt sich ein Temperaturkoeffizient fur die Resonanzfrequenz (wiederum samt ¨
statistischem Fehler). Kommentieren Sie das Ergebnis: Wie kann man das Vorzeichen der Steigung
qualitativ (physikalisch) verstehen?


Kritische Diskussion


------

4.9 Fragen zur Auswertung
• Wozu macht man Fits gemessener Daten?
• Wie lautet der zentrale Grenzwertsatz? Am besten in einem Statistik-Lehrbuch nachschlagen!
• Bei einer Reihe von N Messungen in einem Messpunkt mit Standardabweichung σ, wie groß
ist der statistische Fehler δ des Mittelwerts?
• Bei mehreren Reihen mit je N Messungen im selben Messpunkt, wie ist die Verteilung der
Mittelwerte?
• Was zeichnet die ”
Least Squares“-Methode aus? Wie ist χ
2 definiert?
• Wie schatzt man quantitativ die Qualit ¨ at eines Fits ab? Was f ¨ ur Schl ¨ usse lassen sich aus der ¨
Abschatzung ziehen? 

\end{document}